\chapter{Gegevensoverdracht en foutopsporing}
\label{cha:transfer-and-debug-tools}

De sleutel tot effectieve cross-platform-ontwikkeling is het hebben van snelle en
eenvoudige middelen om programmatuur op de MEGA65 te testen. Dit
geldt zeker nu de MEGA65-emulator nog in ontwikkeling is.
Zelfs wanneer de MEGA65-emulator voltooid is, zal het onwaarschijnlijk
zijn dat deze volledige compatibiliteit op volledige snelheid kan bieden,
omdat de MEGA65 veel veeleisender is om te emuleren dan de C64.

Er zijn verschillende hulpprogramma's die kunnen worden gebruikt voor gegevensoverdracht en
foute opsporing. Deze maken doorgaans gebruik van ofwel de seriële
monitorinterface van de MEGA65, of via de fast-ethernetadapter van de MEGA65. De
seriële monitorinterface is beschikbaar via de UART-lijnen op de JB1-
connector.

Als u geen toegang hebt tot de seriële monitorinterface, worden er
hulpprogramma's ontwikkeld voor de fast-ethernetpoort die een deel,
maar niet alles, van de mogelijkheden van de seriële monitor
bieden. Deze worden gedocumenteerd zodra ze beschikbaar zijn. Het
restant van dit hoofdstuk richt zich op methoden die de seriële
monitorinterface gebruiken.

U kunt een 3.3V UART-adapter op de juiste
lijnen aansluiten, of gemakkelijker, een TE-0790-03 JTAG-debugmodule
op deze connector plaatsen. Dit geeft u een USB-verbinding die kan worden gebruikt
voor het overzetten van software, op afstand fouten opsporen en geheugen inspecteren, alsmede
het activeren of flashen van bitstreams. Met deze verbinding
zijn de volgende hulpprogramma's beschikbaar:

\section{m65-hulpprogramma}

Het repository \url{https://github.com/mega65/mega65-tools} bevat een
aantal hulpprogramma's, utilities en voorbeeldprogramma's. Deze hulpprogramma's zijn voornamelijk voor Linux maar kunnen worden gebruikt op Windows met Cygwin. Een van die hulpprogramma's is
het \stw{m65}-hulpprogramma. Dit is eigenlijk een zakmes
verzameling van utilities in één. Veelgebruikte functies
zijn onder meer:

\subsection{Schermfoto's met het m65-hulpprogramma}

Om een schermfoto van de MEGA65 te maken:

\begin{screencode}
m65 -S
\end{screencode}

Dit maakt een bestand aan met de naam \stw{mega65-screen-000000.png},
of als dit bestand al bestaat, wordt het eerste niet-gebruikte nummer gebruikt
in plaats van \stw{000000}.

Merk op dat deze schermfotofunctie werkt door \stw{m65} de
functionaliteit van de VIC-IV te laten emuleren. Hoewel het uitstekende digitale schermfoto's oplevert, komt het mogelijk niet exact overeen met het werkelijke scherm van de
MEGA65. Op het moment van schrijven rendert het geen sprites of bitplanelen, alleen tekst- en bitmap-gebaseerde videomodi.

\subsection{Een programma laden en uitvoeren op de MEGA65}

Om een programma op de MEGA65 te laden en uit te voeren, kunt u een commando als het volgende gebruiken:

\begin{screencode}
m65 -F -4 -r foo.prg
\end{screencode}

De optie \stw{-F} draagt \stw{m65} op de MEGA65 te herstarten
voordat het programma wordt geladen.

De optie \stw{-4} draagt \stw{m65} op de MEGA65
naar C64-modus te schakelen voordat het programma wordt geladen. Als dit wordt weggelaten,
zal het proberen het programma in C65-modus te laden.

De optie \stw{-r} draagt \stw{m65} op het programma onmiddellijk
na het laden uit te voeren.

Merk op dat dit commando werkt met het normale BASIC LOAD-commando en
daarom beperkt is tot het laden van programma's in het onderste 64 KB van het geheugen.

\subsection{De FPGA herconfigureren om een andere bitstream uit te voeren}

Om een andere MEGA65-bitstream uit te proberen, kan een commando als het volgende worden
gebruikt:

\begin{screencode}
m65 -b bitstream.bit
\end{screencode}

Dit zal de genoemde bitstream naar de FPGA zenden. Omdat de
FPGA door deze actie wordt herconfigureerd, zal het momenteel
uitgevoerde programma niet alleen worden gestopt, maar ook het hoofdgeheugen zal worden
gewist. Voor modellen van de MEGA65 die zijn uitgerust met 8 MB of 16 MB
uitbreidingsgeheugen, worden deze uitbreidingsgeheugens geïmplementeerd in externe
chips en dus zullen de inhoud daarvan niet worden gewist.

Voor niet-MEGA65-bitstreams (zoals zxunomega65 en gbc4mega65), gebruik in plaats daarvan het argument '-q':

\begin{screencode}
m65 -q bitstream.bit
\end{screencode}

\subsection{Toetsenbordinvoer op afstand}

De toetsenbordinterface van de MEGA65 ondersteunt het overzetten van
synthetische toetsgebeurtenissen via de registers \$D615 -- \$D617.
Het \stw{m65}-hulpprogramma gebruikt dit om typen op afstand op de MEGA65
mogelijk te maken op een manier die transparant is voor software. Er zijn drie manieren om dit
te gebruiken:

\begin{screencode}
m65 -t ditiswattekst
\end{screencode}

Deze vorm typt de opgegeven tekst, in dit geval {\em ditiswattekst}, maar
simuleert niet het indrukken van \specialkey{RETURN}. Als u
het indrukken van \specialkey{RETURN} wilt simuleren, gebruikt u \stw{-T}
in plaats van \stw{-t}, bijvoorbeeld:

\begin{screencode}
m65 -T list
\end{screencode}

Dit zou het LIST-commando typen en uitvoeren.

Tot slot is het mogelijk om algemene toetsenbordbesturing op afstand te starten via:

\begin{screencode}
m65 -t -
\end{screencode}

In deze modus wordt elke toets, die op het toetsenbord van de computer
waarop \stw{m65} draait, doorgegeven aan de MEGA65. Merk op dat
niet alle speciale toetsen worden ondersteund en dat er enige vertraging is, dus
het gebruik van de toetsherhalingsfunctie kan onverwachte resultaten opleveren. Voor besturing op afstand in het algemeen is het echter een zeer nuttige voorziening.

\subsection{Unit-testing en logboekondersteuning}

Het \stw{m65}-hulpprogramma ondersteunt unit-testing
direct op MEGA65-hardware. Wanneer de \textit{unit testing-modus} actief is,
wacht \stw{m65} tot de MEGA65 bepaalde byte-reeksen via de
seriële interface stuurt die de huidige status (gestart, geslaagd, mislukt) signaleren
van een gegeven test. Daarnaast is het mogelijk om logboekberichten te verzenden vanaf
de MEGA65 naar de hostcomputer.

De unittestingmodus wordt betreden door \stw{m65} aan te roepen met de vlag \stw{-u}.
Om een BASIC-programma op afstand in C65-modus uit te voeren en tegelijk \stw{m65}
in de \textit{unittestingmodus} te zetten, kan het volgende commando worden gebruikt:

\begin{screencode}
    m65 -Fur attic-ram.prg -w tests.log
\end{screencode}

De opties \stw{-F} en \stw{-r} vertellen \stw{m65} om de MEGA65 te herstarten
voordat het programma "attic-ram.prg" wordt geladen en automatisch uitgevoerd.
De optie \stw{-u} vertelt \stw{m65} vervolgens om in de \textit{unittestingmodus} te gaan
in plaats van te stoppen na het starten van het programma.
De optionele optie \stw{-w} zorgt ervoor dat \stw{m65} de testresultaten aan het
bestand "test.log" toevoegt (het bestand aanmaakt als het nog niet bestaat).

Merk op dat \stw{m65} automatisch de \textit{unittestingmodus} verlaat
als er gedurende meer dan 10 seconden geen teststatussignalen zijn ontvangen.

Er is ondersteuning voor het verzenden van unittestingsignalen naar de hostcomputer vanuit
C- en BASIC 65-programma's:

\subsubsection{Unit-tests gebruiken met C}

De MEGA65 libc bevat ondersteuning voor unit-testing via functies gedefinieerd in
\stw{tests.h} en \stw{tests.c}.

Om het begin van een test te signaleren, includeert u \stw{tests.h} en gebruikt u:

\begin{verbatim}
    unit_test_setup("testName",issueNumber);
\end{verbatim}

waarbij \textit{"testName"} een leesbare naam van de test is (bijvoorbeeld
"VIC-II") en \textit{issueNumber} een verwijzing naar het bijbehorende
bug-issue (bijvoorbeeld het issue-nummer van GitHub).

Na het starten van een test kunnen geslaagde tests worden gesignaleerd met
de functie unit\_test\_ok():

\begin{verbatim}
    unit_test_ok();
\end{verbatim}

Een mislukte test wordt gesignaleerd met unit\_test\_fail():

\begin{verbatim}
    unit_test_fail("fail message");
\end{verbatim}

Elke keer dat de functies unit\_test\_ok() of unit\_test\_fail()
worden aangeroepen, wordt het \textit{sub-issue} van de test (gerapporteerd op de
hostcomputer) verhoogd. Dit maakt het eenvoudiger om meerdere tests in één bestand te combineren en
te identificeren.

U kunt willekeurige logboekberichten verzenden via unit\_test\_log():

\begin{verbatim}
    unit_test_log("hello world from mega65!");
\end{verbatim}

...en ten slotte, wanneer alles klaar is, wordt het einde van unit-testing gesignaleerd
door het gebruik van:

\begin{verbatim}
    unit_test_done();
\end{verbatim}


\subsubsection{Unit-tests gebruiken met BASIC 65}

\stw{b65support.bin} is een machinetaalmodule die ondersteuning biedt voor
unit-testing vanuit BASIC 65. Deze staat in het mega65-tools-
repository op \url{https://github.com/MEGA65/mega65-tools} en kan worden gebouwd met:

\begin{verbatim}
    make bin65/b65support.bin
\end{verbatim}

Het bestand \stw{b65support.bin} is dan beschikbaar in de map \stw{bin65}.

Deze module werkt door de USR-vector om te leiden om de functies
uit te voeren die nodig zijn om te communiceren met de test-host.

In een geautomatiseerd testsenario wilt u mogelijk het binaire bestand \stw{b65support.bin}
overzetten naar het MEGA65-geheugen door middel van \stw{m65}:

    \begin{screencode}
        m65 -@ mega65-tools/bin65/b65support.bin@15fe
    \end{screencode}


Het is natuurlijk ook mogelijk om \stw{b65support.bin} rechtstreeks vanaf de
MEGA65 te laden door de \stw{M65UTILS.D81}-image van de Freezer te koppelen en het volgende commando te geven:

\begin{screencode}
    BLOAD "B65SUPPORT.BIN"
\end{screencode}

Na het laden wordt \stw{b65support.bin} geïnitialiseerd met:

\begin{screencode}
    SYS $1600
\end{screencode}

Na initialisatie biedt \stw{b65support.bin} de volgende functies:

\begin{screencode}
    A=USR(<issueNum>)
\end{screencode}
bereidt een nieuwe test voor met nummer <issueNum> en herstelt het subissue-nummer naar 0.

\begin{screencode}
    A=USR("=<testName>")
\end{screencode}
stelt de testnaam in en stuurt het testsignaal; bijvoorbeeld: \stw{A=USR("=VIC-III")}
stelt de testnaam in op 'VIC-III' en signaleert de hostcomputer dat de test
is gestart.

\begin{screencode}
    A=USR("/<logMessage>)
\end{screencode}
stuurt een logboekbericht naar de hostcomputer.

\begin{screencode}
    A=USR("P")
\end{screencode}
stuurt het signaal 'geslaagd' naar de hostcomputer en verhoogt het subissue-nummer.

\begin{screencode}
    A=USR("F")
\end{screencode}
stuurt het signaal 'test mislukt' naar de hostcomputer en verhoogt
het subissue-nummer.

\begin{screencode}
    A=USR("D")
\end{screencode}
stuurt het signaal 'test voltooid' naar de hostcomputer.

Alle aanroepen retourneren het huidige subissue-nummer of \stw{?ILLEGAL QUANTITY ERROR} in
geval van een ongeldig commando.

\subsubsection{BASIC 65-voorbeeld}

Het volgende is een volledig BASIC 65-voorbeeld dat laat zien hoe de
unittestingfuncties van \stw{m65} worden gebruikt:

\begin{screencode}
    100 rem attic ram cache test
    110 poke $bfffff2,$e0         : rem schakel cache van zolder-ram in
    120 sys $1600                 : rem initialiseer testmodule
    130 a=usr(379)                : rem zet issue-nummer
    140 a=usr("=attic-ram-cache") : rem zet testnaam
    150 bank128:poke0,65          : rem uit voorzorg
    160 b0=$8000000 : b1=$8000100 : rem zolden-ram-gebieden om te testen
    170 for r=0 to $ff
    180   poke b0+r,0             : rem vul gebied 1 met 0
    190   poke b1+r,$ff           : rem vul gebied 2 met $ff
    200 next r
    210 for t=1 to 10             : rem 10 pogingen
    220   poke b0,32              : rem schrijf naar b0
    230   for x=0 to $ff:t1=b1+x
    240     a=peek(b0)            : rem lees uit b0
    250     b=peek(t1):b=peek(t1) : rem lees tweemaal uit t1
    260     ifb<>255 thenf=t:t=11:x=256   : rem dit mag niet gebeuren
    270   next x
    280 next t
    290 if f=0 then begin
    300   print "geen fouten gevonden na";t;"pogingen."
    310   a=usr("p")              : rem signal 'test passed' to host
    320 bend : else begin
    330   a=usr("f")              : rem signal 'test failed' to host
    340   print "hyper-ram-fout gedetecteerd na";f;"pogingen."
    350   print "peek($";hex$(t1);") [t1] is";b;"maar moet 255 zijn"
    360 bend
    370 a=usr("d")                : rem test gereed    
\end{screencode}

\section{M65Connect}

Dit is een cross-platform grafisch hulpprogramma dat beschikbaar is voor Windows, Linux
en MacOSX, waarmee toegang wordt geboden tot de meeste functies van
het \stw{m65}-hulpprogramma, zonder dat een
commandoregel nodig is, of zonder dat u het hulpprogramma voor uw favoriete
besturingssysteem hoeft te compileren.

Het repository voor M65Connect is: \url{https://github.com/MEGA65/m65connect}

De nieuwste binaire versie is beschikbaar op \url{https://files.mega65.org}.

Met de MEGA65 of Nexys FPGA uitgeschakeld, sluit een USB-kabel aan van uw computer op de MEGA65 of Nexys FPGA-board. Voer het \textit{M65Connect}-programma uit en volg de aanwijzingen om verbinding te maken. Het programma helpt u bij het identificeren van welke USB Serial Port moet worden gebruikt.

Met dit hulpprogramma kunt u eenvoudig PRG-programma's en diverse andere bestanden overdragen. M65Connect kan de overdracht, het schakelen naar C64-modus en de uitvoering van programma's afhandelen.

\section{mega65\_ftp}

Het \stw{mega65}\_\stw{ftp}-hulpprogramma van
het repository \url{https://github.com/mega65/mega65-tools} is
enigszins misleidend genoemd: hoewel het
een File Transfer Program is, gebruikt het niet het File Transfer
Protocol (FTP). Het gebruikt in plaats daarvan de seriële monitorinterface om
de MEGA65 op afstand te besturen en heeft directe toegang tot de SD-kaart om het
kopiëren van bestanden tussen de MEGA65 en de hostcomputer mogelijk te maken.

Merk op dat het de staat van de MEGA65 niet perfect herstelt bij het afsluiten,
en dus alleen moet worden gebruikt wanneer de MEGA65 bij de READY.-prompt staat,
zodat eventuele draaiende software niet in de war raakt. Met name moet u
het gebruik vermijden wanneer een gevoelig programma draait, zoals
het Freeze-menu, het MEGA65-configuratiehulpprogramma of het MEGA65
Format/FDISK-hulpprogramma. (Dit probleem kan met weinig
inspanning worden opgelost, als iemand de tijd en interesse heeft om het te fixen).

Bij uitvoering biedt het een FTP-achtige interface die de
commando's \stw{get}, \stw{put}, \stw{rename} en \stw{dir} ondersteunt.
Merk op dat u bij het plaatsen van een bestand moet zorgen dat het een
naam heeft die volledig in hoofdletters staat en een DOS-compatibele 8.3-teken
bestandsnaam heeft. Dit komt door beperkingen in zowel \stw{mega65}\_\stw{ftp} als de
VFAT32-bestandssysteemcode van de Hypervisor van de MEGA65. Opnieuw kunnen deze problemen
worden opgelost met een bescheiden hoeveelheid inspanning van de kant van een
gemotiveerd lid van de gemeenschap.

Tot slot is het programma \stw{mega65}\_\stw{ftp} {\em zeer} langzaam bij het pushen
van nieuwe bestanden naar de MEGA65, met doorgaans snelheden van rond de 5 KB/sec.
Dit komt deels doordat de seriële interface in staat is om
gegevens over te dragen met slechts 40 KB/sec (bij 4.000.000 bits per
seconde), en deels doordat het besturingsproces op afstand resulteert in
veel heen-en-weerverkeer waarbij hulpoutines op de MEGA65 worden uitgevoerd om
sectoren op de SD-kaart te lezen, schrijven en verifiëren. Het zou heel
haalbaar zijn om dit te verbeteren tot bijna 40 KB/sec, en mogelijk
sneller met een combinatie van gegevenscompressie,
deduplicatie van identieke sectoren (met name bij het uploaden van schijf-
images) en andere technieken. Opnieuw zou dit een zeer welkome bijdrage zijn die iemand uit de gemeenschap zou kunnen leveren
ten algemene nutte.

\section{TFTP-server}

Er is een echte TFTP-server voor de MEGA65 in ontwikkeling die snelle
TFTP-overdrachten via de 100 Mbit-ethernet ondersteunt en die kan worden gebruikt om
zeer snel bestanden van de MEGA65 te lezen. Snelheden van bijna 1 MB/sec
zijn mogelijk, afhankelijk van de prestaties van de SD-kaart. In plaats van
DHCP te gebruiken, zal dit hulpprogramma reageren op {\em elk} IP-adres dat eindigt op
.65. Het gebruikt altijd het MAC-adres 40:40:40:40:40:40. Echte DHCP-
ondersteuning en het gebruik van het geconfigureerde ethernet-MAC-adres van de MEGA65
worden mogelijk in de toekomst toegevoegd.

Belangrijker nog, de ondersteuning voor het schrijven van
bestanden naar de SD-kaart is nog niet volledig en wordt geblokkeerd door de behoefte
aan de implementatie van de nodige functies in de Hypervisor van de MEGA65 voor het aanmaken en
vergrootten van bestanden. Een bijzondere uitdaging is het mogelijk maken van het aanmaken van
bestanden met aaneengesloten clusters zoals vereist voor D81-schijfimages: als
een D81-bestand gefragmenteerd is, kan het niet worden gekoppeld, omdat het
koppelingsmechanisme een pointer naar het aaneengesloten blok van de
SD-kaart vereist dat de schijfimage bevat.
In de tussentijd kan \stw{mega65}\_\stw{ftp} als alternatief worden gebruikt.

\section{Een BASIC-tekstbestandlisting omzetten naar een PRG-bestand}

Als u een ongetokeniseerd BASIC-programma in platte-tekstformaat hebt dat afkomstig is van bijvoorbeeld een internetbericht, en u wilt het op de MEGA65 proberen zonder het in te typen, dan is het mogelijk het om te zetten naar een PRG.

\textit{C64List} is een Windows-gebaseerd commandoregelhulpprogramma waarmee u de conversie kunt uitvoeren. Zodra u een \textit{.PRG}-bestand hebt, kunt u een hulpprogramma zoals M65Connect gebruiken om het naar de MEGA65 of Nexys FPGA te uploaden.

C64List is beschikbaar voor download op \url{http://www.commodoreserver.com/Downloads.asp}

Zorg ervoor dat u een programmalijst in een bestand op uw lokale computer hebt opgeslagen (bijvoorbeeld \textit{program.txt}) gecodeerd als ANSI of UTF8.

Gebruik C64List om het bestand om te zetten naar een PRG-bestand:
    
\begin{tcolorbox}[colback=black,coltext=white]
    \begin{verbatim}
        C64List program.txt -prg
    \end{verbatim}
\end{tcolorbox}

Nu kunt u uw nieuw geconverteerde programma uploaden naar de MEGA65 met M65Connect of een van de eerder beschreven hulpprogramma's.

Het is vermeldenswaard dat deze methode niet 100\% effectief zal zijn bij listings met speciale PETSCII-tekens. Programma's met PETSCII vereisen enige bewerking op de MEGA65 zelf voordat ze op schijf worden opgeslagen.
